\chapter{Background and Related Work}\label{sec:literature-review}

\section{Classical Relay Strategies}\label{sec:classical-relay-strategies}

The fundamental question that motivates this thesis is: can a learned relay function \(f_\theta(\cdot)\) complement classical relay functions (AF scaling, DF regeneration) when their channel assumptions are mismatched or the relevant mapping is difficult to specify analytically? On the matched memoryless channel, the relevant symbol statistics are already captured by the classical MAP rule; the potential advantage of learning must therefore be established empirically in low-SNR or mismatched regimes rather than presumed.

\subsection{Amplify-and-Forward (AF)}\label{sec:classical-relay-strategies-af}

The AF relay amplifies the received signal by a gain factor \(G\) that normalizes the output power: (Equation~\ref{eq:af-gain})
\addcontentsline{equ}{listequations}{\protect\numberline{\theequation}\quad af-gain}

\begin{equation}\protect\phantomsection\label{eq:af-gain}{G = \sqrt{\frac{P_{\text{target}}}{\mathbb{E}[|y_R|^2]}}
}\end{equation}

For CSI-assisted variable-gain AF, the standard dual-hop end-to-end SNR is \cite{Laneman2004CooperativeDiversity,TseViswanath2005Fundamentals}: (Equation~\ref{eq:af-snr})

\begin{equation}\protect\phantomsection\label{eq:af-snr}{\text{SNR}_{\text{eff}}^{\text{AF}} = \frac{\text{SNR}_1 \cdot \text{SNR}_2}{\text{SNR}_1 + \text{SNR}_2 + 1}
}\end{equation}
\addcontentsline{equ}{listequations}{\protect\numberline{\theequation}\quad af-snr}

This form assumes \emph{CSI-assisted variable-gain} AF (Laneman et al.), in which the relay computes a fresh gain from its instantaneous first-hop channel state. That is different from the implementation in this thesis: its gain is a single block-adaptive scalar estimated from the empirical received power over the simulated block. It is therefore data- and noise-dependent and requires the block to be available before the gain is fixed; it is neither conventional statistical fixed-gain AF nor per-symbol CSI-assisted variable-gain AF. Equations~\ref{eq:af-snr}--\ref{eq:af-p-eff} are included only as standard background and are not closed forms for the simulated AF results, which come entirely from Monte Carlo evaluation.

This reveals AF's fundamental limitation: the effective SNR is bottlenecked by the weaker hop even when the other has high SNR, because the relay amplifies the first hop's noise along with its signal, accumulating both at the destination.

\subsection{Decode-and-Forward (DF)}\label{sec:classical-relay-strategies-df}
The DF relay demodulates the received signal, recovers the transmitted bits, and re-modulates clean symbols:

\begin{equation}\protect\phantomsection\label{eq:df-demod-remod}{\hat{b}_R = \text{demod}(y_R), \quad x_R = \text{mod}(\hat{b}_R)
}\end{equation}
\addcontentsline{equ}{listequations}{\protect\numberline{\theequation}\quad df-demod-remod}

The end-to-end BER for DF is: (Equation~\ref{eq:df-ber-hops})

\begin{equation}\protect\phantomsection\label{eq:df-ber-hops}{P_e^{\text{DF}} = P_{e,1}(1 - P_{e,2}) + (1 - P_{e,1}) \cdot P_{e,2}
}\end{equation}
\addcontentsline{equ}{listequations}{\protect\numberline{\theequation}\quad df-ber-hops}

where \(P_{e,1}\) and \(P_{e,2}\) are the BER of the first and second hops, respectively: an end-to-end bit error occurs exactly when \emph{one} of the two hops flips the bit (two flips cancel), so this expands to \(P_{e,1} + P_{e,2} - 2P_{e,1}P_{e,2}\), the same expression used in the AWGN analysis (Equation~\ref{eq:df-ber-awgn}). For BPSK modulation over AWGN (real-valued noise with variance \(1/(2\gamma)\)), each hop's BER is \(P_e = Q\left(\sqrt{2\gamma}\right)\) on the simulator's \(\gamma=E_b/N_0\) axis. DF provides clean regeneration at high SNR but suffers from error propagation when the first-hop BER is non-negligible. As discussed in Remark~\ref{rem:df-terminology}, the ``DF'' baseline used here is the uncoded symbol-wise slicer, not practical finite-length coded block-DF; its optimality claims are therefore confined to the uncoded setting.

\subsection{Other Classical Relay Techniques}

AF and DF are two canonical relay strategies: AF linearly forwards an analog observation, whereas DF regenerates a detected symbol or decoded block. Compress-forward, estimate-forward, and compute-forward address different network models and objectives and are not intermediate points on a single AF--DF scale. The focus on AF and symbol-wise DF is deliberate but limited: they are representative relay-function baselines for the canonical benchmark, not strict performance bounds or an exhaustive account of classical relaying.

\section{Machine Learning in Wireless Communication}\label{sec:machine-learning-in-wireless-communication}

Applying machine learning to physical-layer wireless communication has gained significant momentum, driven by neural networks' ability to learn complex non-linear mappings directly from data where analytical solutions are intractable or suboptimal: channel estimation \cite{YeLiJuang2018OFDMDL}, signal detection \cite{SamuelDiskinWunder2019MIMODetect}, autoencoder-based end-to-end communication \cite{Dorner2018OverTheAirDL}, and resource allocation \cite{Sun2018LearnToOptimize}.

\subsection{Theoretical Basis: Universal Approximation and Denoising}\label{sec:theoretical-basis-universal-approximation-and-denoising}

The theoretical justification for applying neural networks to relay signal processing rests on two pillars. First, \textbf{universal approximation} results \cite{Cybenko1989UAT,Hornik1991UAT,GoodfellowBengioCourville2016DeepLearning} show that feedforward networks with suitable non-polynomial activations can approximate continuous functions on compact domains to arbitrary accuracy, given sufficient width. This existence result does not guarantee finite-sample learning, optimization, or generalization. For relay denoising, the target function maps noisy observations to clean transmitted symbols:

\begin{equation}\protect\phantomsection\label{eq:optimal-denoiser}{f^*: \mathbb{R}^{2w+1} \to [-1, 1], \quad f^*(y_{i-w}, \dots, y_{i+w}) = \mathbb{E}[x_i \mid y_{i-w}, \dots, y_{i+w}]
}\end{equation}
\addcontentsline{equ}{listequations}{\protect\numberline{\theequation}\quad optimal-denoiser}

This conditional expectation \(f^*\) is the Bayes-optimal denoiser: it minimizes the mean squared error (MSE) over all possible estimators. For BPSK symbols corrupted by AWGN, \(f^*\) reduces to the posterior mean:

\begin{equation}\protect\phantomsection\label{eq:optimal-denoiser-awgn}{f^*(y) = \tanh\left(\frac{y}{\sigma^2}\right)
}\end{equation}
\addcontentsline{equ}{listequations}{\protect\numberline{\theequation}\quad optimal-denoiser-awgn}

which is a smooth sigmoid approaching the hard-decision signum function as \(\sigma^2 \to 0\). A single tanh unit with the right input scaling represents this posterior \emph{exactly} for a fixed \(\sigma^2\). The deployed relay, however, is trained at SNR \(\in\{5,10,15\}\)~dB and evaluated over \(0\)--\(20\)~dB, so one set of weights must approximate a whole \emph{family} of posteriors indexed by \(\sigma^2\), and on the fading channel a mixture over random effective SNRs. That a 169-parameter network approximates that family closely enough to track the classical baselines is an empirical finding of this thesis, not a consequence of the exact single-\(\sigma^2\) representation.

Second, the \textbf{bias-variance decomposition} provides a framework for understanding model complexity:

\begin{equation}\protect\phantomsection\label{eq:bias-variance-decomp}{\mathbb{E}[(\hat{x} - x)^2] = \text{Bias}^2(\hat{x}) + \text{Var}(\hat{x}) + \sigma^2_{\text{irreducible}}
}\end{equation}
\addcontentsline{equ}{listequations}{\protect\numberline{\theequation}\quad bias-variance-decomp}

For the denoising problem considered here the irreducible term is the Bayes risk $\mathbb{E}[\operatorname{Var}(X \mid Y)]$, the conditional uncertainty that remains after observing $Y$, rather than the channel-noise variance $\sigma^2$ itself. The decomposition is used as a conceptual framework; the irreducible term should be read that way. Increasing capacity can change both bias and variance, but it does not necessarily reduce one or increase the other in modern regularized networks. Since the target \(f^*\) is a smooth low-dimensional mapping in the memoryless special case, the working hypothesis is that additional capacity may buy little. Whether that occurs is an empirical question; Section~\ref{sec:capacity-and-overfitting} reports only what was measured and does not infer overfitting without train--test-gap evidence.

\subsection{Prior Work in Deep Learning for Physical-Layer Processing}\label{sec:prior-work-in-deep-learning-for-physical-layer-processing}

Ye et al.~\cite{YeLiJuang2018OFDMDL} demonstrated joint DNN-based channel estimation and signal detection in OFDM systems using received pilot and data information. Their approach replaces separately implemented estimation and equalization stages and was reported to be robust under several modeled impairments; it does not eliminate pilots.

Dorner et al.~\cite{Dorner2018OverTheAirDL} proposed treating the modulator, channel, and demodulator as an autoencoder trained end-to-end with a differentiable message-reconstruction objective and evaluated by error rate. This approach jointly optimizes the communication chain rather than directly minimizing the nondifferentiable BER.

The sampled relay-assisted literature includes learning for channel estimation and network-level decisions rather than only the forwarding map. Akdemir et al.~\cite{Akdemir2024DLRelaySelection} use learned channel-coefficient forecasts within conventional DF systems and compare SNR-based, threshold-based and all-relay participation protocols. This is not evidence that a neural relay-selection classifier directly replaces the relay's signal processing. No systematic search establishes the prevalence of either approach across the field.

Model-based learned detectors place neural components inside classical recursions: Tsai et al.~\cite{TsaiTeng2020BCJRNet} learn channel-dependent likelihoods for neural BCJR detection, while Karanov et al.~\cite{Karanov2024BurstyISI} study ISI with bursty noise. Chen et al.~\cite{Chen2024BCJRNetRobustness} report benefits under imperfect channel knowledge, with robustness depending on channel variation. These are examples of the importance of model match, not proof of a universal boundary. Section~\ref{sec:bcjr-benchmark} supplies a QPSK relay-output BCJR control; it is not an end-to-end learned-relay comparison.

A broader perspective on machine learning's role across wireless protocol layers, including cooperative communications, is given in \cite{Gunduz2019MLAir}. A separate line of work exploits the structural similarity between AF relay networks and neural networks, treating relay nodes as neurons and tuning relay gains with deep-learning optimization to exploit hardware nonlinearities \cite{Bergel2023RelaysAsNeurons,BergelICASSP2024}; the gains reported there are on analog gain optimization, not symbol-level detection.

Classical AF/DF frameworks remain well understood and widely adopted \cite{Nosratinia2004CooperativeCommunication} but rely on fixed linear forwarding or hard-decision rules. Learning the relay's symbol-level input-output mapping itself has drawn comparatively little attention, motivating the data-driven alternative studied here.

\subsection{Neural Network Relay Processing}\label{sec:neural-network-relay-processing}

For relay processing specifically, a supervised learning approach trains a neural network \(f_\theta\) to minimize the empirical mean-squared error: (Equation~\ref{eq:mse-loss})

\begin{equation}\protect\phantomsection\label{eq:mse-loss}{\mathcal{L}(\theta) = \frac{1}{N} \sum_{i=1}^{N} (\hat{x}_i - x_i)^2, \quad \hat{x}_i = f_\theta(y_{i-w:i+w})
}\end{equation}
\addcontentsline{equ}{listequations}{\protect\numberline{\theequation}\quad mse-loss}

where \(\hat{x}_i = f_\theta(y_{i-w:i+w})\) is the network output based on a sliding window of \(2w+1\) received symbols, and \(x_i\) is the clean transmitted symbol. This window-based approach provides temporal context that enables the network to exploit statistical dependencies in the noise-corrupted signal.

A well-known example of \(f_\theta(\cdot)\) is the \emph{multilayer perceptron (MLP)}, a feedforward artificial neural network composed of an input layer, one or more fully connected hidden layers with nonlinear activation functions, and an output layer, trained using backpropagation to model nonlinear input--output relationships.
The MLP equation is shown in equation \ref{eq:mlp-relay}:

\begin{equation}\protect\phantomsection\label{eq:mlp-relay}{\mathbf{h} = \text{ReLU}(\mathbf{W}_1 \mathbf{w} + \mathbf{b}_1), \quad \hat{x} = \tanh(\mathbf{W}_2 \mathbf{h} + \mathbf{b}_2)
}\end{equation}
\addcontentsline{equ}{listequations}{\protect\numberline{\theequation}\quad mlp-relay}

The sliding window formulation ($i \in [i-w:i+w]$) is retained for architectural consistency with the unknown-ISI experiments of Chapter~\ref{sec:unknown-channels}. Under the canonical i.i.d.\ per-symbol fading model, adjacent observations carry no additional information about the current symbol, and the window is redundant there by construction.
Under the canonical i.i.d.\ assumptions, neighboring observations do not add information about the current symbol conditional on its channel state. For hard BPSK/QPSK slicing, phase-corrected equalization preserves the decision boundary. A soft estimator also depends on the instantaneous reliability: the equalized sample alone is not generally sufficient for the CSI-conditioned posterior when the channel magnitude is discarded. The unknown-channel study separately examines temporal context under ISI.

\textbf{Multi-SNR training} is a key design choice: training on data generated at multiple SNR levels (5, 10, 15 dB here) yields a single estimator whose denoising function generalizes across operating conditions rather than specializing to one noise level. In decision-theoretic terms the training objective averages the loss over the sampled SNRs, so what it minimises is \emph{Bayes} risk under the empirical SNR distribution, not the worst-case (minimax) risk over that range: nothing in the objective protects the SNR at which the network happens to do worst.

For resource-limited relays, it is useful to test whether a modest model suffices for the chosen processing task. Universal approximation and bias--variance reasoning do not predict a mandatory BER upturn as neural-network width increases. The capacity sweep therefore treats plateauing or degradation as empirical hypotheses, not theoretical consequences.

\section{Generative Models for Signal Processing}\label{sec:generative-models-for-signal-processing}

Generative models offer an alternative to directly learning a denoising function (the discriminative approach above): they instead learn the distribution of clean signals \(p(\mathbf{x})\). Bayesian denoising provides the general probabilistic framing:
\begin{equation}\protect\phantomsection\label{eq:bayes-rule}{p(\mathbf{x} | \mathbf{y}) \propto p(\mathbf{y} | \mathbf{x}) p(\mathbf{x})}
\end{equation}

The generative approach is theoretically appealing because it separates the modeling of the signal prior from the noise model, potentially enabling better generalization to unseen noise conditions. The VAE implemented in this thesis does not compute that posterior explicitly. It is a latent-variable generative architecture trained supervised, with noisy observations at the encoder and the clean symbol as the reconstruction target, under a KL regulariser; Eq.~\eqref{eq:bayes-rule} frames the problem it addresses rather than describing the operation it performs.

\subsection{Variational Autoencoders}\label{sec:variational-autoencoders}

\textbf{Variational Autoencoders (VAEs)} \cite{KingmaWelling2014VAE} learn a latent representation by maximizing the evidence lower bound (ELBO). The generative model posits that data \(\mathbf{x}\) is generated from a latent variable \(\mathbf{z} \sim p(\mathbf{z}) = \mathcal{N}(\mathbf{0}, \mathbf{I})\) through a decoder \(p_\theta(\mathbf{x} | \mathbf{z})\). Since the true posterior \(p(\mathbf{z} | \mathbf{x})\) is intractable, an encoder network \(q_\phi(\mathbf{z} | \mathbf{x})\) approximates it. The ELBO objective is derived from the log-evidence decomposition: (Equation~\ref{eq:vae-elbo})

\begin{equation}\protect\phantomsection\label{eq:vae-elbo}{\log p_\theta(\mathbf{x}) = \underbrace{\mathbb{E}_{q_\phi}\left[\log \frac{p_\theta(\mathbf{x}, \mathbf{z})}{q_\phi(\mathbf{z} | \mathbf{x})}\right]}_{\text{ELBO}(\theta, \phi; \mathbf{x})} + \underbrace{D_{KL}(q_\phi(\mathbf{z} | \mathbf{x}) \| p_\theta(\mathbf{z} | \mathbf{x}))}_{\geq 0}
}\end{equation}
\addcontentsline{equ}{listequations}{\protect\numberline{\theequation}\quad vae-elbo}

Since the KL divergence is non-negative, the ELBO is a lower bound on the log-evidence. Maximizing the ELBO simultaneously trains the encoder to approximate the true posterior and the decoder to reconstruct the data. The ELBO decomposes into a reconstruction term and a regularization term:

\begin{equation}\protect\phantomsection\label{eq:vae-loss}{\mathcal{L}_{\text{VAE}} = \underbrace{\mathbb{E}_{q_\phi(\mathbf{z}|\mathbf{x})}[\log p_\theta(\mathbf{x}|\mathbf{z})]}_{\text{Reconstruction quality}} - \underbrace{\beta \cdot D_{KL}(q_\phi(\mathbf{z}|\mathbf{x}) \| p(\mathbf{z}))}_{\text{Latent space regularity}}
}\end{equation}
\addcontentsline{equ}{listequations}{\protect\numberline{\theequation}\quad vae-loss}

The reconstruction term encourages accurate recovery while the KL term regularizes the latent distribution toward the prior. The $\beta$ weight in Eq.~\eqref{eq:vae-loss} changes that tradeoff \cite{Higgins2017betaVAE}. Choosing $\beta<1$ puts more relative weight on reconstruction; it does not by itself establish a BER benefit for this relay task. Larger $\beta$ motivates more constrained representations, but disentanglement is not guaranteed or measured here.

The \emph{reparameterization-trick} enables gradient-based optimization through the stochastic sampling layer:
\begin{equation}\protect\phantomsection\label{eq:repremtrized-trick}
\mathbf{z} \sim q_\phi(\mathbf{z} | \mathbf{x}) = \mathcal{N}(\boldsymbol{\mu}, \text{diag}(\boldsymbol{\sigma}^2))
\end{equation}
Instead of sampling from equation \ref{eq:repremtrized-trick}, the sample is expressed as:
\begin{equation}\protect\phantomsection\label{eq:repremtrized-trick-resample}
\mathbf{z} = \boldsymbol{\mu} + \boldsymbol{\sigma} \odot \boldsymbol{\epsilon}
\end{equation}
\addcontentsline{equ}{listequations}{\protect\numberline{\theequation}\quad repremtrized-trick}

where \(\boldsymbol{\epsilon} \sim \mathcal{N}(\mathbf{0}, \mathbf{I})\).
This moves the stochasticity outside the computational graph, allowing backpropagation through the encoder.

The implemented relay is a supervised denoising VAE rather than the unconditional model above: its encoder receives the noisy observation window, while reconstruction is evaluated against the clean target. The latent regularization supplies a generative inductive bias, but the denoising property follows from this noisy-input/clean-target training objective rather than from the standard VAE derivation alone. Sampling \(\mathbf{z}\) at inference also injects variance into an otherwise deterministic estimation task; Section~\ref{sec:generative-vs.-discriminative-paradigms-for-relay-processing} reports the measured outcome without attributing a causal effect to that sampling.

\subsection{Generative vs.~Discriminative Paradigms for Relay Processing}\label{sec:generative-vs.-discriminative-paradigms-for-relay-processing}

The application of generative models to relay processing has, to the best of our knowledge, not been extensively studied. The key question is whether the generative inductive bias (learning the data distribution rather than just the input-output mapping) provides any advantage for the relay denoising task. For BPSK, the clean signal distribution is trivially simple (a discrete distribution on \(\{-1, +1\}\)), suggesting that the generative overhead may not be justified. This thesis compares VAE-based relay processing against classical and supervised learning methods on a common benchmark, and the results show that the generative paradigm provides no significant advantage for this particular task. A conditional GAN relay is implemented and described in Chapter~\ref{sec:methods} but is excluded from the reported comparison on cost grounds, so no claim is made about it here. On the canonical setup the VAE tracks the supervised MLP and symbol-wise DF rather than trailing them, reaching $0.00972$ at 20~dB against DF's $0.00972$ and the MLP's $0.00992$ (Table~\ref{tbl:table2}), and it stays inside the feedforward group at every SNR in the equal-parameter comparison as well (Table~\ref{tbl:table8}). The generative bias is therefore neither an advantage nor a handicap here: it costs parameters and training time to arrive where a 169-parameter regressor already is. One qualification remains, unchanged by this: the relay samples $\mathbf{z}\sim q_\phi(\mathbf{z}\mid\mathbf{x})$ stochastically at inference, and the deterministic variant that substitutes the posterior mean $\boldsymbol{\mu}$ at test time was not evaluated, so nothing here bears on whether that variant would do better.

\section{Sequence Models: Transformers, State Space Models}\label{sec:sequence-models-transformers-and-state-space-models}

Two sequence-modelling paradigms are evaluated as relay functions in this thesis, and what matters here is
their inductive bias and their cost rather than their internal machinery, which is standard and is not
reproduced. \textbf{Transformers} \cite{Vaswani2017Attention} use multi-head self-attention. In the bidirectional encoder configuration used here, every position
attends to every other through learned query, key and value projections, giving a global receptive field over the buffered window; a causal (decoder) Transformer instead attends only to the past, and would not see the right-hand half of the symmetric window this thesis feeds its relays. The cost is
$O(n^2)$ time and memory in the window length $n$. Because attention is permutation-equivariant, sinusoidal
positional encodings supply the ordering the relay needs to distinguish symbols within its window. For the
11-symbol window used throughout, the quadratic term is negligible; what the architecture costs is
parameters, four projection matrices per head.

\textbf{Structured state space models} \cite{GuGoelRe2022SSM} take the opposite view, treating the sequence
model as a discretised linear dynamical system --- a learned recursive filter, linear in the sequence length,
with the HiPPO initialisation supplying a principled memory of the past. \textbf{Mamba} \cite{gu2023mamba}
makes that system time-varying by conditioning its parameters on the input, so the state-transition term acts
as a learned gate that can reset or retain the state per symbol; \textbf{Mamba-2} \cite{DaoGu2024TransformersSSM}
recasts a structured selective SSM through its duality with a structured attention matrix. SSD uses a more
restrictive scalar-identity transition structure than a general diagonal SSM; that restriction enables the
faster implementation, so Mamba-2 is not simply the same function class with a faster kernel. The classical analogue is adaptive filtering, and that
connection becomes relevant only in the unknown-channel study of Chapter~\ref{sec:unknown-channels}; on the
canonical memoryless channel there is no temporal structure for a recurrence to exploit.

\subsection{State Space Models for Signal Processing}\label{sec:state-space-models-for-signal-processing}

The application of state space models as the \emph{relay-side} processing function is, to the best of our
knowledge, largely unexamined. Transformers and state-space models have already been applied to a range of
wireless receiver and signal-processing tasks; what the literature reviewed here did not turn up is a
controlled comparison of these architectures against one another, at a matched parameter budget, in the
two-hop relay benchmark studied in this thesis. Within that reviewed literature, no prior controlled
comparison of Mamba-S6, Mamba-2 SSD and Transformers for the same relay-processing task was identified.

\section{Theoretical Foundations: Channel Models}\label{sec:theoretical-foundations-channel-models}

This section presents the theoretical BER analysis for each channel model used in this study, derives the closed-form expressions, and validates them against Monte Carlo simulation. The theoretical-simulative comparison serves two purposes: (i) it verifies the correctness of the simulation framework, and (ii) it establishes the baseline performance that AI relays must beat. This matters beyond internal bookkeeping, but the coverage has to be stated exactly, because it is partial. The closed forms derived here validate the memoryless components to which they apply --- the AWGN and Rayleigh symbol-error expressions, the two-hop DF composition built from them, and the channel and noise conventions every experiment shares. They do \emph{not} validate every classical baseline used later. The AF relay implemented in this thesis normalizes by a block-adaptive empirical-power scalar rather than an instantaneous CSI-assisted gain, so Eq.~\eqref{eq:af-snr} is background for it and not a prediction of its curve (Section~\ref{sec:classical-relay-strategies-af}). Viterbi MLSE, the BCJR/APP detector and CMA have no closed-form BER on the channels they are run on, and are audited instead by implementation-specific controls stated where they appear: that exact ML decoding returns an admissible path with likelihood at least as large as that of the transmitted sequence, that it reduces to the nearest-symbol slicer when the ISI is removed, and that the BER-optimal rule agrees with symbol-MAP where the two provably coincide. The comparisons that follow therefore set a learned relay against classical detectors verified in these two different ways, by closed form where one exists and by control otherwise, rather than against unaudited simulations of them.

\subsection{AWGN Channel: Theoretical Analysis}\label{sec:awgn-channel-theoretical-analysis}

The AWGN channel adds zero-mean Gaussian noise to the transmitted signal: (Equation~\ref{eq:awgn-channel})

\begin{equation}\protect\phantomsection\label{eq:awgn-channel}{y = x + n, \quad n \sim \mathcal{N}(0, \sigma^2), \quad \sigma^2 = \frac{P_s}{2\,\gamma}
}\end{equation}
\addcontentsline{equ}{listequations}{\protect\numberline{\theequation}\quad awgn-channel}

where \(P_s = \mathbb{E}[|x|^2] = 1\) for unit-power BPSK symbols and \(\gamma=10^{\mathrm{SNR}_{\mathrm{dB}}/10}=E_s/N_0\). The factor of two is the standard \(N_0/2\) variance per real dimension.

\subsubsection{Single-hop BER}\label{sec:awgn-single-hop-ber}
 For BPSK \(\pm 1\) symbols and real-valued AWGN with variance \(\sigma^2 = 1/(2\gamma)\), an error occurs when the noise pushes the received sample past the decision boundary at the origin. The theoretical BER is: (Equation~\ref{eq:awgn-ber})

\begin{equation}\protect\phantomsection\label{eq:awgn-ber}{P_e^{\text{AWGN}} = Q\!\left(\frac{1}{\sigma}\right) = Q\left(\sqrt{2\gamma}\right) = \frac{1}{2}\,\text{erfc}\!\left(\sqrt{\gamma}\right)
}\end{equation}
\addcontentsline{equ}{listequations}{\protect\numberline{\theequation}\quad awgn-ber}

where \(Q(x) = \frac{1}{2}\text{erfc}(x/\sqrt{2})\) is the Gaussian Q-function. The simulator's SNR label is the normalized \(E_s/N_0\) axis, not the raw ratio \(P_s/\sigma^2\) in a single real dimension.

\begin{remark}[Noise convention]\label{noise-convension}
The standard textbook result \(P_e = Q(\sqrt{2 E_b/N_0})\) assumes a one-sided noise PSD \(N_0\), hence variance \(N_0/2\) per real dimension. The implementation uses \(\sigma^2=P_s/(2\gamma)\), so for BPSK (\(E_s=E_b=P_s\)) its BER is \(Q(\sqrt{2\gamma})\), exactly the textbook expression with \(\gamma=E_b/N_0\). For QPSK, \(E_b=E_s/2\), so the same simulator label corresponds to an \(E_b/N_0\) value 3.01~dB lower.

For complex baseband signals the same convention gives total noise power \(N_0=P_s/\gamma\), split equally between I and Q; after coherent equalization each real decision dimension therefore has variance \(P_s/(2\gamma|h|^2)\). The Rayleigh expressions use that per-dimension variance and consequently retain the standard factor of two in the Q-function argument.
\end{remark}

\subsubsection{Two-hop AF BER}\label{sec:awgn-two-hop-af-ber}
For the analytical CSI-assisted variable-gain AF model, let the per-hop SNRs be symbol energy divided by complex noise power. With equal per-hop SNR $\gamma$, the effective symbol SNR is:

\begin{equation}\protect\phantomsection\label{eq:af-snr-eff}{\text{SNR}_{\text{eff}}^{\text{AF}} = \frac{\text{SNR}_1 \cdot \text{SNR}_2}{\text{SNR}_1 + \text{SNR}_2 + 1} = \frac{\gamma^2}{2\gamma + 1}
}\end{equation}
\addcontentsline{equ}{listequations}{\protect\numberline{\theequation}\quad af-snr-eff}

where
\(\gamma = \text{SNR}\) is the per-hop SNR. The resulting BER is

\begin{equation}\protect\phantomsection\label{eq:af-p-eff}{P_e^{\text{AF,BPSK}} = Q\left(\sqrt{2\,\text{SNR}_{\text{eff}}}\right)
}\end{equation}
\addcontentsline{equ}{listequations}{\protect\numberline{\theequation}\quad af-p-eff}

for coherent BPSK, using the $N_0/2$ noise variance per decision dimension (Remark~\ref{noise-convension}). For Gray-QPSK the per-bit expression is $Q(\sqrt{\text{SNR}_{\text{eff}}})$. At high SNR,

\begin{equation}\protect\phantomsection\label{eq:gamma-high-snr}{\text{SNR}_{\text{eff}} \approx \frac{\gamma}{2}}
\end{equation}
\addcontentsline{equ}{listequations}{\protect\numberline{\theequation}\quad gamma-high-snr}

confirming the well-known 3 dB penalty of AF relaying. This penalty is a property of the analytical CSI-assisted variable-gain model above. The AF relay implemented in this thesis uses a single block-adaptive scalar estimated from the empirical received power, so Eq.~\eqref{eq:gamma-high-snr} is background rather than a closed-form prediction of the simulated AF curve, and the 3 dB figure should not be read as explaining it.

\subsubsection{Two-hop DF BER}\label{sec:awgn-two-hop-df-ber}
For a decode-and-forward relay with equal-SNR hops, an error occurs at the destination if exactly one hop introduces an error:

\begin{equation}\protect\phantomsection\label{eq:df-ber-awgn}{P_e^{\text{DF}} = P_{e,1} + P_{e,2} - 2P_{e,1}P_{e,2}
}\end{equation}
\addcontentsline{equ}{listequations}{\protect\numberline{\theequation}\quad df-ber-awgn}

With equal hops (i.e.\ \(P_{e,1} = P_{e,2} = P_e^{\text{AWGN}}\)), this becomes

\begin{equation}\protect\phantomsection\label{eq:equal-hop-snr}{P_e^{\text{DF}} = 2P_e(1 - P_e)}
\end{equation}\addcontentsline{equ}{listequations}{\protect\numberline{\theequation}\quad equal-hop-snr}

At high SNR, \(P_e \ll 1\) and
\begin{equation}\protect\phantomsection\label{eq:df-high-snr-two-hop-penalty}{P_e^{\text{DF}} \approx 2P_e}
\end{equation}\addcontentsline{equ}{listequations}{\protect\numberline{\theequation}\quad df-high-snr-two-hop-penalty}
i.e., the two-hop penalty is approximately a factor of 2 in BER.

\subsection{Rayleigh Fading Channel: Theoretical Analysis}\label{sec:rayleigh-fading-channel-theoretical-analysis}

The Rayleigh fading channel models non-line-of-sight (NLOS) propagation where the signal undergoes multiplicative fading:

\begin{equation}\protect\phantomsection\label{eq:rayleigh-channel}{y = hx + n, \quad h \sim \mathcal{CN}(0, 1), \quad n \sim \mathcal{CN}(0, \sigma^2)
}\end{equation}
\addcontentsline{equ}{listequations}{\protect\numberline{\theequation}\quad rayleigh-channel}

The fading coefficient \(h\) is a circularly-symmetric complex Gaussian random variable, so its magnitude \(|h|\) follows a Rayleigh distribution:

\begin{equation}\protect\phantomsection\label{eq:rayleigh-pdf}{f_{|h|}(r) = 2r \cdot e^{-r^2}, \quad r \geq 0
}\end{equation}
\addcontentsline{equ}{listequations}{\protect\numberline{\theequation}\quad rayleigh-pdf}

with \(\mathbb{E}[|h|^2] = 1\) (unit average power). The instantaneous SNR after equalization (\(\hat{x} = y/h\)) becomes \(\gamma_h = |h|^2 \cdot \text{SNR}\), which is exponentially distributed.

\subsubsection{Single-hop BER}\label{sec:rayleigh-single-hop-ber}
Averaging the conditional BER \(P_e(\gamma_h) = Q(\sqrt{2\gamma_h})\) over the exponential distribution of \(\gamma_h\) yields the closed-form \cite[Eq.~14-4-15]{ProakisSalehi2008DigitalComm} (Equation~\ref{eq:rayleigh-ber}):

\begin{equation}\protect\phantomsection\label{eq:rayleigh-ber}{P_e^{\text{Rayleigh}} = \frac{1}{2}\left(1 - \sqrt{\frac{\bar{\gamma}}{1 + \bar{\gamma}}}\right)
}\end{equation}
\addcontentsline{equ}{listequations}{\protect\numberline{\theequation}\quad rayleigh-ber}

where \(\bar{\gamma} = \gamma_b = \bar{E_b}/N_0\) is the average SNR \emph{per bit}. This distinction is load-bearing, because the canonical experiments plot \(\gamma_s = E_s/N_0\) per \emph{symbol}: for QPSK, \(\gamma_b = \gamma_s / 2\), a \(3.01\)~dB offset. Substituting gives the per-bit Rayleigh expression in terms of the plotted axis, \(P_b = \tfrac{1}{2}\bigl(1 - \sqrt{\gamma_s/(\gamma_s+2)}\bigr)\), which is the form the DF-theory column of Chapter~\ref{sec:experiments} evaluates. At high SNR (\(\bar{\gamma} \gg 1\)):

\begin{equation}\protect\phantomsection\label{eq:rayleigh-ber-approx}{P_e^{\text{Rayleigh}} \approx \frac{1}{4\bar{\gamma}}
}\end{equation}
\addcontentsline{equ}{listequations}{\protect\numberline{\theequation}\quad rayleigh-ber-approx}

This \(1/\text{SNR}\) decay is fundamentally slower than the exponential decay of AWGN (\(Q(\sqrt{2\gamma}) \sim e^{-\gamma}/\sqrt{4\pi\gamma}\)), explaining why Rayleigh fading is significantly more challenging. The channel is \textbf{diversity-limited}: deep fades (where \(|h| \approx 0\)) cause errors regardless of the average SNR, and a single link of this kind has diversity order one. Cooperative relaying is often introduced as the remedy, by giving the destination a second, independently faded look at the same message. That is not the mechanism at work here, and it should not be read as the motivation for the architectures this thesis studies. The canonical setup has \emph{no} direct source--destination path (Table~\ref{tbl:assumptions-claims}), so the source's message reaches the destination along one route only and no combining is possible; the two-hop DF expression derived immediately below, \(P_e^{\text{DF}} = 2P_e(1 - P_e)\), is itself first order in \(P_e\) and therefore describes a diversity order of one, not two. What relaying buys on such a link is a shorter hop and hence a higher per-hop SNR at fixed transmit power, at the cost of a second error opportunity. The question this thesis asks is a different one again: given that geometry, what the relay should \emph{compute} on the signal it receives.

\subsubsection{Two-hop DF BER}\label{sec:rayleigh-two-hop-df-ber}
The two-hop DF relay over Rayleigh fading follows the same composition rule as AWGN:

\begin{equation}\protect\phantomsection\label{eq:df-ber-rayleigh}{P_e^{\text{DF,Rayleigh}} = 2P_e^{\text{Rayleigh}}(1 - P_e^{\text{Rayleigh}})
}\end{equation}
\addcontentsline{equ}{listequations}{\protect\numberline{\theequation}\quad df-ber-rayleigh}

\section{Channel Coding: Convolutional Codes and Trellis Decoding}\label{sec:channel-coding-convolutional-codes}

The channel models above assume an uncoded link. Chapter~\ref{sec:experiments}'s supplementary coded study
(Section~\ref{sec:coded-block-df}) departs from that, so this section fixes the three quantities that study
reports on. The underlying coding theory is standard and is not reproduced here; see Proakis and
Salehi~\cite{ProakisSalehi2008DigitalComm} for convolutional encoding and the Viterbi
algorithm~\cite{Forney1972MLSE}, and Bahl \emph{et al.}~\cite{BahlCockeJelinekRaviv1974BCJR} for the
forward--backward (BCJR) recursion. Two properties of those two decoders matter later and are used without
re-derivation: Viterbi returns the maximum-likelihood \emph{sequence} while BCJR returns per-symbol
posteriors, so the two optimise different criteria and neither dominates the other on every metric
(Section~\ref{sec:bcjr-benchmark} measures this distinction); for a binary rate-$1/n$ code, a constraint-length-$K$
shift register has a $2^{K-1}$-state trellis. The principal coded study uses $k=1$, $n=2$, $K=3$ and obtains
rates $2/3$ and $3/4$ from that mother code by puncturing. A separate code-strength control also evaluates $K=5$ and $K=7$.

A code's rate $R = k/n$ and the constellation size $M$ together set the spectral efficiency

\begin{equation}\protect\phantomsection\label{eq:spectral-efficiency}{\eta = R \cdot \log_2 M \quad \text{[bits/symbol]}}\end{equation}
\addcontentsline{equ}{listequations}{\protect\numberline{\theequation}\quad spectral-efficiency}

so a rate-$1/2$ code halves the information rate of whatever modulation carries it, and comparing a coded
link against an uncoded one at the same SNR without accounting for this is not a like-for-like comparison
(Section~\ref{sec:coded-latency-throughput} revisits Chapter~\ref{sec:experiments}'s coded-versus-uncoded
comparison on that basis). Spectral efficiency counts bits transmitted rather than bits that arrived, so the
quantity a link actually adapts on is the goodput

\begin{equation}\protect\phantomsection\label{eq:goodput}{G = \eta \cdot (1 - \mathrm{FER}) \quad \text{[information bits/symbol]}}\end{equation}
\addcontentsline{equ}{listequations}{\protect\numberline{\theequation}\quad goodput}

Maximizing $G$ over a grid of modulation-and-coding schemes at each SNR --- the standard adaptive modulation
and coding (AMC) problem --- traces the link-adaptation envelope that Section~\ref{sec:coded-rate-adaptation}
measures.
